\documentclass[12pt]{article}

%\usepackage{showkeys}
\usepackage{color}
\usepackage{amsmath}
\usepackage{amssymb}
\definecolor{blue}{rgb}{0,0,1}
%%\usepackage[dvips]{graphics}
\usepackage{graphicx}
\usepackage{cite}
\usepackage{mathtools}

\newcommand{\cb}{\textcolor{red}}
\newcommand{\as}{\textcolor{blue}}

\newtheorem{theorem}{Theorem}
\newtheorem{lemma}{Lemma}
\newtheorem{corol}{Corollary}
\newtheorem{prop}{Property}
\newtheorem{propo}{Proposition}
\newtheorem{rem}{Remark}
\newtheorem{dfn}{Definition}
\newtheorem{example}{Example}
\newcommand{\fin}{\,\rule{1ex}{1.6ex}\,}

\date{~}

\textheight 22.6cm \textwidth 17.7cm \evensidemargin -1cm
\oddsidemargin -1cm \topmargin -1cm


\title{The life and the work of Andr\'{e} Louis Cholesky}
\author{Claude Brezinski\thanks{Universit\'{e} de Lille, CNRS, UMR 8524 - Laboratoire Paul Painlev\'{e},
F-59000 Lille, France. E-mail: {\tt Claude.Brezinski@univ-lille.fr}.}
}


\begin{document}

\maketitle

In this talk, I first describe the life of Andr\'{e} Louis Choleky (1873-1918) who was a French army officer specialised in topography and cartography. Then, I analyse his scientific work. In particular, I discuss his well known method for solving a system of linear equations with a symmetric positive definite matrix. I show how this method was forgotten and then came back to light. The other works of Cholesky will also be mentioned.

\end{document}